\documentclass[a4paper]{article}
\usepackage{modulITERA}
\begin{document}
\SampulDok{Kontrak}{Kontrak Kuliah}{Kesepakatan antara tim pengajar dan mahasiswa tentang tujuan, jadwal, penilaian, dan aturan Praktikum Pemrograman.}
\par\vspace{4pt}\noindent{\fontsize{9.5pt}{13pt}\selectfont
\begin{tabular}{@{}>{\raggedright\arraybackslash}p{0.280\textwidth}>{\raggedright\arraybackslash}p{0.700\textwidth}@{}}
\kepala{Komponen} & \kepala{Keterangan}\\
Kode dan nama & IF25-11002 Praktikum Pemrograman\\\garis
Bobot & 2 SKS praktikum (16 pertemuan × 150 menit)\\\garis
Semester & 1 (Ganjil)\\\garis
Prasyarat & tidak ada\\\garis
Mata kuliah pasangan & IF25-11001 Algoritma Pemrograman, topik sejajar setiap minggu\\\garis
Program studi & S1 Teknik Informatika, Fakultas Teknologi Industri, Institut Teknologi Sumatera\\\garis
Tim pengajar & Tim Pengajar Praktikum Pemrograman\\\garis
Bahasa dan alat & C++ (standar C++17); GDB Online di kelas, g++ di komputer sendiri\\\garis
Tahun & 2026\\
\garisdasar
\end{tabular}}\par\vspace{6pt}
\newpage
\Bagian{Kontrak}{Manfaat dan deskripsi}
Praktikum Pemrograman melatih mahasiswa menerjemahkan algoritma menjadi program C++ yang berjalan, lalu menjelaskan mengapa program itu berjalan. Materi mencakup struktur program, tipe data dan masukan, operator, percabangan, perulangan, array satu dan dua dimensi, string, pencarian, pengurutan, fungsi, rekursi, dan struct. Setiap pertemuan reguler memakai siklus problem-first lima fase: review, struggle, materi, hands-on, dan Exit Ticket.

Mata kuliah ini menjadi fondasi untuk hampir semua mata kuliah pemrograman berikutnya. Yang dilatih bukan hanya menulis kode, tetapi juga menelusuri dan menjelaskan kode, karena kemampuan itulah yang dipakai saat membaca dan memperbaiki program orang lain.

\Bagian{Kontrak}{Capaian pembelajaran}
\par\vspace{4pt}\noindent{\fontsize{9.5pt}{13pt}\selectfont
\begin{tabular}{@{}>{\raggedright\arraybackslash}p{0.160\textwidth}>{\raggedright\arraybackslash}p{0.820\textwidth}@{}}
\kepala{Kode} & \kepala{Rumusan}\\
CPL06 & Mampu menganalisis dan mengimplementasi konsep dasar matematika, statistika, algoritma dan sistem komputer sebagai fondasi komputasional.\\\garis
CPMK0607 & Mahasiswa mampu menerapkan konsep dasar algoritma dan sistem komputer untuk menyelesaikan masalah komputasi. Diukur lewat Bagian A setiap instrumen.\\\garis
CPMK0608 & Mahasiswa mampu menjelaskan konsep dasar algoritma dan sistem komputer untuk menyelesaikan masalah komputasi. Diukur lewat Bagian B setiap instrumen.\\
\garisdasar
\end{tabular}}\par\vspace{6pt}
\Bagian{Kontrak}{Jadwal dan tugas}
\par\vspace{4pt}\noindent{\fontsize{9pt}{12pt}\selectfont
\begin{longtable}{@{}>{\raggedright\arraybackslash}p{0.060\textwidth}>{\raggedright\arraybackslash}p{0.480\textwidth}>{\raggedright\arraybackslash}p{0.440\textwidth}@{}}
\kepala{P} & \kepala{Topik} & \kepala{Tugas}\\
\endhead
1 & Hello World dan Struktur Program C++ & Tugas 1: Kartu Perkenalan Digital\\\garis
2 & Variabel, Tipe Data, Ekspresi, dan Input & Tugas 2: Program Pendaftaran Mahasiswa Baru\\\garis
3 & Operator dan Precedence & Tugas 3: Kalkulator Lengkap dengan Validasi\\\garis
4 & Percabangan & Tugas 4: Sistem Tiket Bioskop\\\garis
5 & Perulangan Bagian 1 & Tugas 5: Program Statistik Nilai\\\garis
6 & Perulangan Bagian 2 dan Pola & Tugas 6: Pattern Generator\\\garis
7 & Review dan Simulasi UTS & tanpa tugas\\\garis
8 & Ujian Tengah Semester & contoh soal dan kisi-kisi\\\garis
9 & Array Satu Dimensi & Tugas 9: Sistem Nilai Mahasiswa\\\garis
10 & Array Dua Dimensi dan String & Tugas 10: Sistem Nilai Kelas\\\garis
11 & Searching & Tugas 11: Direktori Kontak\\\garis
12 & Sorting & Tugas 12: Sistem Leaderboard\\\garis
13 & Fungsi dan Modularitas & Tugas 13: Library Matematika\\\garis
14 & Rekursi dan Struct & Tugas 14: Perpustakaan Mini dengan Struct dan Rekursi\\\garis
15 & Review dan Simulasi UAS & tanpa tugas\\\garis
16 & Ujian Akhir Semester & contoh soal dan kisi-kisi\\
\garisdasar
\end{longtable}}\par\vspace{6pt}
\Bagian{Kontrak}{Penilaian}
\par\vspace{4pt}\noindent{\fontsize{9.5pt}{13pt}\selectfont
\begin{tabular}{@{}>{\raggedright\arraybackslash}p{0.150\textwidth}>{\raggedright\arraybackslash}p{0.380\textwidth}>{\raggedright\arraybackslash}p{0.120\textwidth}>{\raggedright\arraybackslash}p{0.160\textwidth}>{\raggedright\arraybackslash}p{0.160\textwidth}@{}}
\kepala{Komponen} & \kepala{Pelaksanaan} & \kepala{Bobot} & \kepala{CPMK0607} & \kepala{CPMK0608}\\
Exit Ticket & 14 kali (P1–P7, P9–P15), rata-rata & 25\% & 12,5 & 12,5\\\garis
Tugas & 12 kali (P1–P6, P9–P14), rata-rata & 25\% & 12,5 & 12,5\\\garis
UTS & P8, 4 soal, 120 menit & 25\% & 12,5 & 12,5\\\garis
UAS & P16, 4 soal, 150 menit & 25\% & 12,5 & 12,5\\\garis
\textbf{Total} &  & \textbf{100\%} & \textbf{50} & \textbf{50}\\
\garisdasar
\end{tabular}}\par\vspace{6pt}
Setiap instrumen dibagi rata antara Bagian A (kode) dan Bagian B (tracing dan penjelasan). Nilai minimum lulus setiap CPMK adalah 50; di bawah itu ada remedial. Konversi nilai huruf mengikuti A-03 Standar Penilaian Pembelajaran ITERA.

\Bagian{Kontrak}{Aturan yang disepakati}
\par\vspace{4pt}\noindent{\fontsize{9.5pt}{13pt}\selectfont
\begin{tabular}{@{}>{\raggedright\arraybackslash}p{0.200\textwidth}>{\raggedright\arraybackslash}p{0.780\textwidth}@{}}
\kepala{Hal} & \kepala{Kesepakatan}\\
Kehadiran & Mengikuti ketentuan kehadiran dalam peraturan akademik ITERA. Exit Ticket hanya bisa dikerjakan oleh mahasiswa yang hadir.\\\garis
Exit Ticket & Dikerjakan di 25 menit terakhir, sendiri, tanpa AI, internet, atau diskusi. Pelanggaran membuat Exit Ticket pertemuan itu bernilai 0.\\\garis
Tugas & Dikumpulkan sebelum pertemuan berikutnya lewat kanal yang diumumkan tim pengajar. Ketentuan keterlambatan diumumkan pada pertemuan pertama.\\\garis
Penggunaan AI & Boleh untuk memahami pesan error atau konsep saat hands-on dan tugas, tidak untuk meminta solusi utuh. Bagian B tugas harus memuat penjelasan dan pengalaman sendiri.\\\garis
Kejujuran akademik & Menyalin pekerjaan teman atau AI pada instrumen yang dinilai membuat nilai instrumen itu 0 dan dapat dilaporkan ke program studi.\\\garis
UTS dan UAS & Boleh membawa satu lembar catatan A4 tulisan tangan sendiri. Soal B dikerjakan di kertas tanpa menjalankan program.\\\garis
Konsultasi & Konsultasi dilakukan langsung dengan tim pengajar.\\
\garisdasar
\end{tabular}}\par\vspace{6pt}
\Bagian{Kontrak}{Persetujuan}
Kontrak ini dibacakan dan disepakati pada Pertemuan 1. Perubahan hanya dilakukan atas kesepakatan tim pengajar dan perwakilan mahasiswa, lalu diumumkan kepada seluruh kelas.

\vspace{14mm}
\noindent\begin{tabular}{@{}p{0.45\textwidth}p{0.45\textwidth}@{}}
Tim Pengajar & Perwakilan Mahasiswa\\[22mm]
\rule{55mm}{0.5pt} & \rule{55mm}{0.5pt}\\
Nama dan tanda tangan & Nama, NIM, dan tanda tangan\\
\end{tabular}
\end{document}
